\documentclass[11pt,a4paper]{article}
\usepackage[utf8]{inputenc}
\usepackage[T1]{fontenc}
\usepackage{geometry}
\usepackage{amsmath,amssymb}
\usepackage{booktabs}
\usepackage{graphicx}
\usepackage[numbers]{natbib}
\usepackage{listings}
\usepackage{tabularx}
\usepackage[colorlinks=true,linkcolor=blue,citecolor=blue,urlcolor=blue]{hyperref}

\geometry{
  left=2.5cm,
  right=2.5cm,
  top=2.5cm,
  bottom=2.5cm
}

\title{\textbf{NRR-Energy: Calibration of Commit/Defer Utility\\Under Existing Paired Inputs}}

\author{
  Kei Saito\thanks{ORCID: 0009-0006-4715-9176} \\
  Independent Researcher, Japan \\
  \texttt{kei.saito.research@gmail.com}
}

\date{March 2026 \\[0.5em]
\textit{Part of the Non-Resolution Reasoning (NRR) research program.}}

\begin{document}
\maketitle

\begin{center}
\textcopyright\ 2026 Kei Saito.
Licensed under CC BY 4.0.\\
\url{https://creativecommons.org/licenses/by/4.0/}
\end{center}

\begin{abstract}
This paper introduces an Energy-layer calibration procedure for commit/defer decision support using existing paired datasets carried forward from integrated paper7 package-first surfaces, including selected Stage B and ordered-combo slices. The method fixes utility construction and calibration-only thresholding before test-side evaluation. We define utility from quality-improvement and rework-cost terms and evaluate two cost formulations: gross rework improvement and output-overhead-adjusted net cost. Under tested conditions, net-cost calibration expands decision coverage relative to gross-cost calibration, but this expansion is condition-sensitive. In Stage B, net-cost calibration increases coverage while also increasing disagreement against calibration-derived oracle labels; in ordered-combo data, net-cost calibration increases coverage without increasing overall disagreement. Sensitivity analyses identify non-effective regions concentrated in Gemini and in Stage B at higher temperature. Narrow implementation-gated forward checks also replay exactly on the tested covered slices while exposing a provider-specific transfer boundary: Anthropic remains route-stable on the tested slices, whereas Gemini loses stable-slice route stability on the tested combo opening. Because realized forward outcomes remain matched between baseline and Energy arms in these batches, these forward checks support boundary mapping rather than gain claims. Additional downstream dialogue-side checks show a similar bounded pattern: early priority control remains reproducible across ambiguity, preacceptance-burden, and repaired first-step side-topic/quarantine surfaces on the tested provider pair, but the main incremental first-step lift over baseline concentrates in \texttt{C\_incident\_response} and later followthrough remains conditional rather than uniformly stable. Taken together, the results establish a calibrated boundary map for Energy instrumentation and delimit where direct policy claims require additional forward paired validation.
\end{abstract}

\noindent\textbf{Keywords:} Non-Resolution Reasoning, Energy calibration, commit/defer control, boundary conditions, disagreement analysis

\section{Introduction}
\label{sec:intro}

NRR-Energy is one module in the ongoing NRR research program:
Core $\rightarrow$ Phi $\rightarrow$ IME $\rightarrow$ Transfer $\rightarrow$ Coupled $\rightarrow$ integrated paper7 $\rightarrow$ Energy.
Core and Phi established state- and entropy-side operator behavior~\cite{saito2025nrr,saito2026phi}. IME and Transfer established extraction reliability and transfer behavior~\cite{saito2026ime,saito2026transfer}. Coupled added dependency-propagation checks~\cite{saito2026coupled}. Integrated paper7 then fixed the delayed-commitment comparison surface together with the selected provider-sensitive boundary-honesty layer that Energy reuses through carried-forward paired inputs~\cite{saito2026principles}. Energy uses those paired artifacts to calibrate a commit/defer utility instrument and to report where that instrument is stable versus non-effective.

This manuscript remains a calibration paper. It answers two main questions under fixed procedures: (i) can an Energy utility instrument be defined and reproduced, and (ii) under which conditions does that calibrated instrument produce stable versus unstable decision behavior. It also reports narrow downstream boundary checks: whether a frozen routed Stage A path and first covered forward slices can execute and replay exactly, and whether later dialogue followthrough remains condition-bounded even after early priority control succeeds, without upgrading the evidence to a policy-improvement claim. It does not answer whether an online Energy policy improves outcomes prospectively versus baseline under new paired runs.

NRR is not an anti-LLM framework.
NRR does not replace standard LLM use.
NRR optimizes when to commit and when to defer, under explicit conditions.
NRR is a derivation-and-evaluation framework, not a closed recipe; the current operator and condition set is an evaluated set under tested conditions, not an exhaustive set.

\textbf{Series Consistency Statement:}
At this evidence state, unsupported interpretations are not injected, hypothesis-set expansion is evidence-gated, and unobserved possibilities remain explicitly unresolved.

\subsection{Contributions}

\begin{enumerate}
  \item A reproducible calibration procedure for an Energy-style commit/defer utility instrument under existing paired inputs carried forward from integrated paper7 package-first surfaces.
  \item Multi-profile calibration reporting across quality- and cost-weighted utility settings.
  \item Explicit boundary reporting for non-effective and disagreement-prone regions, including provider, temperature, and ordered-combination effects.
  \item A clear claim boundary that separates calibration evidence, narrow downstream boundary checks, implementation-gated forward boundary checks, and future gain validation.
\end{enumerate}

\section{Framework and Calibration Procedure}
\label{sec:framework}

\subsection{Inputs}
\label{sec:inputs}

The calibration uses three input classes:
\begin{itemize}
  \item Stage B paired metrics (\texttt{stats/stageb\_all/mst\_pair\_metrics.csv})
  \item ordered-combo paired metrics (\texttt{stats/combo\_rep3\_all/mst\_pair\_metrics.csv})
  \item raw run traces for output-token overhead reconstruction
\end{itemize}

The current split instantiation uses calibration rows from reps 1 and 2 and test rows from rep 3. This split instance is provisional; the split procedure itself is the fixed part.

\subsection{Utility Definition}
\label{sec:utility}

We define a quality term and two cost terms:
\begin{itemize}
  \item $Q =$ misinterpretation improvement
  \item $C_{\mathrm{gross}} =$ rework-cost improvement
  \item $C_{\mathrm{net}}(\mathrm{output}) = C_{\mathrm{gross}} - \mathrm{overhead}_{\mathrm{output}}$
\end{itemize}

The main calibration utility is
\[
\Delta U = w_Q \cdot z(Q)_{\mathrm{cal}} + w_C \cdot z(C)_{\mathrm{cal}},
\]
where z-scoring is computed on calibration rows only.
Here $C$ denotes either $C_{\mathrm{gross}}$ or $C_{\mathrm{net}}(\mathrm{output})$ depending on the calibration variant.
The gross variant serves as a reference calibration and the net variant as the primary calibration reported in Sections~\ref{sec:results} and~\ref{sec:boundaries}.

\subsection{Profiles, Threshold, and Oracle}
\label{sec:profiles}

Three weight profiles are reported:
\begin{itemize}
  \item \texttt{q80\_c20}
  \item \texttt{q50\_c50}
  \item \texttt{q20\_c80}
\end{itemize}

For each condition and profile, the tie threshold is derived as
\[
\tau = q_{0.2}\left(|\Delta U_{\mathrm{cal}}|\right).
\]

The oracle sign per condition is defined by the calibration-side mean sign of $\Delta U$. Test rows then emit a decision (\texttt{prefer\_pattern}, \texttt{prefer\_baseline}, or \texttt{tie}), a coverage flag, and a disagreement flag. Reported disagreement rates use Wilson 95\% confidence intervals~\cite{wilson1927} from $(\mathrm{disagree\_count}, n_{\mathrm{test}})$.

\subsection{E-1 Freeze Map}

\begin{table}[ht]
\centering
\small
\begin{tabularx}{\textwidth}{@{}lXl@{}}
\toprule
E-1 item & Freeze decision in this draft & Method anchor \\
\midrule
E1-01 & Adopted (utility frame). Semantic interpretation is finalized jointly with E1-08 cost definition. & Secs.~\ref{sec:utility}, \ref{sec:costgate} \\
E1-02 & Adopted (calibration-only z-score normalization). & Sec.~\ref{sec:utility} \\
E1-03 & Adopted (multi-profile sensitivity; three weight profiles). & Secs.~\ref{sec:profiles}, \ref{sec:e1summary} \\
E1-04 & Adopted (quantile threshold derivation). & Sec.~\ref{sec:profiles} \\
E1-05 & Adopted as procedure; current rep-based split instance is provisional. & Sec.~\ref{sec:inputs} \\
E1-06 & Adopted (CSV-complete outputs for tie, coverage, disagreement, and CI-ready counts). & Sec.~\ref{sec:profiles} \\
E1-07 & Adopted (noise-floor method retained as diagnostic only, not as the primary decision rule). & Sec.~\ref{sec:profiles} \\
E1-08 & Adopted (\texttt{C\_net(output)} primary, \texttt{C\_gross} reference, \texttt{C\_net(total)} rejected for primary use). & Sec.~\ref{sec:costgate} \\
E1-09 & Adopted as procedure freeze (derive $\tau$ per condition/profile), not value freeze of a single global table. & Secs.~\ref{sec:profiles}, \ref{sec:e1summary} \\
\bottomrule
\end{tabularx}
\caption{Mapping from the E-1 worksheet to the current calibration procedure.}
\label{tab:e1_map}
\end{table}

\section{Calibration Results}
\label{sec:results}

\subsection{E-1 Calibration Summary}
\label{sec:e1summary}

\begin{table}[ht]
\centering
\small
\begin{tabular}{@{}llrrrr@{}}
\toprule
Dataset & Profile & $\tau$ (gross) & $\tau$ (net) & Coverage gross $\rightarrow$ net & Disagreement gross $\rightarrow$ net \\
\midrule
Stage B & \texttt{q80\_c20} & 0.01231 & 0.01184 & 57.30\% $\rightarrow$ 74.16\% & 8.99\% $\rightarrow$ 14.61\% \\
Stage B & \texttt{q50\_c50} & 0.01428 & 0.01377 & 52.81\% $\rightarrow$ 84.27\% & 6.74\% $\rightarrow$ 17.98\% \\
Stage B & \texttt{q20\_c80} & 0.04087 & 0.03349 & 51.69\% $\rightarrow$ 80.90\% & 6.74\% $\rightarrow$ 11.24\% \\
Combo & \texttt{q80\_c20} & 0.08366 & 0.08101 & 47.83\% $\rightarrow$ 71.01\% & 2.90\% $\rightarrow$ 2.90\% \\
Combo & \texttt{q50\_c50} & 0.04310 & 0.03941 & 47.83\% $\rightarrow$ 75.36\% & 2.90\% $\rightarrow$ 2.90\% \\
Combo & \texttt{q20\_c80} & 0.00255 & 0.02245 & 63.77\% $\rightarrow$ 75.36\% & 7.25\% $\rightarrow$ 4.35\% \\
\bottomrule
\end{tabular}
\caption{Calibration summary across all three weight profiles.}
\label{tab:e1_summary}
\end{table}

Interpretation follows four stable points. First, net calibration increases coverage in all dataset-profile cells. Second, Stage B shows coverage gain with disagreement increase under all profiles. Third, Combo shows coverage gain with stable-or-lower disagreement. Fourth, in Combo the \texttt{q20\_c80} profile changes the gross-side disagreement level (7.25\%) relative to the other profiles (2.90\%), so stronger cost weighting alters the gross calibration behavior itself rather than only the net correction. Across profiles, Stage B remains higher-disagreement than Combo under net calibration.

\subsection{Cost Definition Gate}
\label{sec:costgate}

Observed sign-flip rates for $C_{\mathrm{gross}} \rightarrow C_{\mathrm{net}}$ are 37.08\% (99/267) for Stage B and 38.21\% (81/212) for Combo under \texttt{C\_net(output)}, versus 66.29\% (177/267) and 65.09\% (138/212) under \texttt{C\_net(total)}. We therefore adopt \texttt{C\_net(output)} as the primary cost term and reject \texttt{C\_net(total)} for primary calibration because input-side fixed cost dominates and distorts recoverable-cost interpretation. \texttt{C\_gross} remains the reference upper-bound view.

\subsection{E-2 Sensitivity}

The primary reporting slice is \texttt{q50\_c50} under net calibration. Overall disagreement is 17.98\% for Stage B and 2.90\% for Combo, with coverage of 84.27\% and 75.36\%, respectively. Provider-separated disagreement under the same slice is 10.00\% (3/30) for Anthropic, 34.48\% (10/29) for Gemini, and 10.00\% (3/30) for OpenAI in Stage B, versus 0.00\% (0/24) for Anthropic, 9.52\% (2/21) for Gemini, and 0.00\% (0/24) for OpenAI in Combo. Temperature sensitivity in Stage B remains ordered as $t=0.3$ (11/45; 24.44\%) $>$ $t=0.0$ (5/44; 11.36\%). In Combo, weaker directions appear in \texttt{conditional\_hierarchical} and \texttt{delta\_hierarchical}, both at 11.11\% disagreement.

Across all three profiles, Gemini remains the highest-disagreement provider in both datasets, Stage B preserves the $t=0.3 > t=0.0$ ordering, and \texttt{conditional\_hierarchical} plus \texttt{delta\_hierarchical} remain weaker cells in Combo.

\subsection{Implementation-Gated Forward Boundary Check}
\label{sec:forwardboundary}

Under the current \texttt{decision\_context\_v2 + compact-grounding} path, routed live Stage A freeze reached \texttt{frozen} for Anthropic and Gemini in both \texttt{stable} and \texttt{context-drift} conditions. This resolves the earlier wiring blocker at the implementation-trial level, but it does not by itself establish that the routed shared-prefix readout is a direct grounded-state-quality measure.

Forward covered-slice checks then produced exact replay on both tested providers for the branch slice (\(24/24\) exact matches each). Anthropic kept the stable scenario on the pattern side (\(8/8\)) while \texttt{context-drift} and \texttt{drift-return-recovery} stayed baseline-side (\(8/8\) each). Gemini completed the same covered slice with exact replay, but its stable scenario split between pattern and baseline (\(3/8\) versus \(5/8\)), establishing a provider-specific route-instability boundary even before combo widening.

The first branch-adjacent two-direction combo opening sharpened that contrast. Anthropic replayed exactly on \(48/48\) pairs and kept the stable slice pattern-side for both tested directions (\(16/16\) total), whereas Gemini also replayed exactly on \(48/48\) pairs but sent the same stable combo slice fully baseline-side (\(16/16\)). In all of these forward batches, realized baseline and Energy outcomes remained matched within batch. The forward evidence is therefore operational and boundary-facing: it shows that the current instrument can be executed and audited on the tested slices, while also delimiting a provider-specific transfer boundary that prevents any gain claim at this stage.

\subsection{Dialogue-Side Carryforward Boundary}
\label{sec:dialoguecarryforward}

The current downstream \texttt{priority-resolution} package provides a bounded dialogue-side carryforward read under the tested provider pair. Early control surfaces were reproducible across three checks: ambiguity resolution without unilateral commitment, downstream preacceptance-burden removal, and a repaired first-step side-topic/quarantine surface. On that repaired first-step surface, the revised Energy policy was clean across the tested rows, but the incremental lift over baseline was not uniform: baseline was already fairly strong on \texttt{A\_trip\_planning} and \texttt{B\_hiring\_packet}, while the main additional lift concentrated in \texttt{C\_incident\_response}. This carryforward read remains bounded. It does not establish provider-universal side-topic robustness, and later-horizon followthrough remains conditional rather than uniformly stable.

\section{Boundary Conditions}
\label{sec:boundaries}

\begin{table}[ht]
\centering
\small
\begin{tabularx}{\textwidth}{@{}Xcc@{}}
\toprule
Region (\texttt{q50\_c50}, net) & Disagree rate & 95\% Wilson CI \\
\midrule
Stage B overall (16/89) & 17.98\% & [11.38\%, 27.23\%] \\
Combo overall (2/69) & 2.90\% & [0.80\%, 9.97\%] \\
Stage B Gemini (10/29) & 34.48\% & [19.94\%, 52.65\%] \\
Combo Gemini (2/21) & 9.52\% & [2.65\%, 28.91\%] \\
Stage B \texttt{t=0.3} (11/45) & 24.44\% & [14.24\%, 38.67\%] \\
Stage B \texttt{t=0.0} (5/44) & 11.36\% & [4.95\%, 23.98\%] \\
Combo \texttt{conditional\_hierarchical} (1/9) & 11.11\% & [1.99\%, 43.50\%] \\
Combo \texttt{delta\_hierarchical} (1/9) & 11.11\% & [1.99\%, 43.50\%] \\
\bottomrule
\end{tabularx}
\caption{Representative non-effective and disagreement-prone regions.}
\label{tab:boundary_regions}
\end{table}

These are boundary findings, not anomalies to hide. Three readings matter operationally. First, Gemini-centered instability is supported, especially in Stage B. Second, Stage B at higher temperature remains a disagreement-risk region. Third, Combo weaker directions are visible but still have wide intervals because $n$ is small; they are boundary warnings rather than high-precision estimates.

\section{Discussion}
\label{sec:discussion}

The Energy instrument can be calibrated reproducibly on existing paired data and can expose where decision confidence is likely to degrade. Calibration quality is nevertheless condition-dependent. In small-effect single-operator settings, output-overhead-adjusted utility can increase callable decisions while also increasing misalignment risk. In larger-effect ordered-combo settings, the same correction improves callable coverage without increasing overall disagreement. Energy should therefore be treated as a condition-gated instrument, not as a universal replacement rule.

The same $\Delta U$ framing can also be read as a bounded cost-of-resolution instrument. Under this read, a commit-side move corresponds to executing deferred comparative resolution at the current output boundary, whereas a defer-side move keeps unresolved alternatives active and postpones that resolution step. A positive $\Delta U$ therefore means not that cost disappears, but that under the tested condition the delayed-resolution route has higher net utility than immediate collapse at that step. The non-effective regions reported here, especially Gemini-centered cells and higher-temperature Stage B cells, then mark conditions where that delay no longer yields a stable net advantage. This remains a local interpretation of the present calibration signal, not a general complexity claim.

In series terms, that makes Energy a calibration and boundary layer rather than the
place where the whole downstream closure is proved. Its role is to separate where
local carryforward looks trustworthy from where later policy claims still require
their own direct tests.

One plausible interpretation is that the Stage B versus Combo disagreement gap reflects effect-size structure rather than a categorical difference between single-operator and ordered-combination datasets. Under tested conditions, Combo cells include stronger directional separations, so the same calibration procedure can expand callable coverage with less disagreement increase than in the smaller-effect Stage B cells. This remains a conditional interpretation: the current evidence supports an effect-size-consistent explanation under these datasets, not a universal law that ordered combinations are always easier to calibrate.

A second boundary is self-reference. The current oracle labels are derived from calibration-split mean $\Delta U$ signs, so the evaluation remains an internal consistency check of the calibrated instrument rather than an external prospective policy test. That is why the next step remains a new paired validation (baseline versus Energy policy) under frozen procedures. Until that forward test is run, the present results justify calibration and boundary reporting, but they do not justify direct claims that the calibrated instrument improves online commit/defer policy outcomes by itself.

A third boundary concerns the routed Stage A implementation path that this calibration is intended to support later. Under current implementation-gated live checks, the routed shared-prefix readout is best treated as a bounded response-quality calibration signal rather than as a direct grounded-state-quality measure. The grounding-sensitive part of the readout is retained by construction and remains relevant for ablation-style checks, but under the currently tested live conditions that part saturates, so the observed live variation is carried mainly by anti-genericity and anti-redundancy terms. We therefore treat the present Stage A path as an implementation-gated calibration path for later forward validation, not as standalone evidence that the routed prefix signal already identifies grounded state quality directly.

A fourth boundary comes from the first covered forward checks. Exact replay and successful routed freeze show that the current implementation path is operationally usable on the tested slices, but operational usability is not the same as balanced forward effectiveness. Anthropic remained route-stable on the tested branch and branch-adjacent combo slices, whereas Gemini lost stable-slice route stability when the combo slice was opened. Because baseline and Energy outcomes remained matched inside those batches, the forward evidence narrows transfer boundaries rather than upgrading this paper to a policy-improvement result.

A fifth boundary comes from the current \texttt{priority-resolution} package that sits downstream of this calibration work. On the tested provider pair, early control surfaces were reproducible: models could resolve ambiguous priority without unilateral commitment, keep preacceptance burden out of the response more reliably, and, under the repaired frozen first-step side-topic/quarantine surface, reach clean first-step side-topic workstart. That side-topic gain was bounded rather than uniform. Baseline was already fairly strong on \texttt{A\_trip\_planning} and \texttt{B\_hiring\_packet}, while the main incremental lift over baseline concentrated in \texttt{C\_incident\_response}. However, later-horizon followthrough remained conditional rather than uniformly stable. Across later checkpoints, both providers often preserved the active-versus-suspended objective boundary while still degrading on unresolved-relation handling, especially by committing too early to an ordering or interpretation that the user had not yet fixed. This indicates that the current Energy path is bounded not only by whether a model can keep the old objective suspended, but also by whether prompt-layer control can preserve unresolved relational structure across later turns.

\section{Conclusion}
\label{sec:conclusion}

NRR-Energy contributes a calibrated commit/defer utility instrument under existing paired inputs carried forward from integrated paper7 and a narrow implementation-gated forward boundary read on how that instrument transfers to frozen routed execution. The current evidence supports reproducible calibration and boundary reporting, not prospective policy-superiority claims. Under tested conditions, net-cost calibration expands callable coverage relative to gross-cost calibration, but the quality of that expansion depends on provider, temperature, and ordered-combination cell. Gemini-centered instability, Stage B at higher temperature, and specific Combo directions remain explicit non-effective regions, and the first forward combo opening adds a further provider-specific transfer boundary on Gemini stable slices. The current downstream \texttt{priority-resolution} package adds a matching dialogue-side boundary: early priority control remains reproducible through ambiguity, downstream preacceptance-burden, and repaired first-step side-topic/quarantine surfaces, but the side-topic lift over baseline is concentrated in \texttt{C\_incident\_response} and later followthrough still weakens at template-dependent depths.

The practical contribution is therefore condition-bounded instrumentation: a way to measure when a commit/defer signal is more or less trustworthy before promoting it to a live policy rule. The first forward checks show that frozen execution and exact replay are achievable on the tested slices, but they do not yet show improvement because baseline and Energy outcomes remained matched within batch. The next step is a broader forward paired validation of baseline versus Energy policy under the frozen procedure defined here, with the current Gemini transfer boundary treated as an explicit watch point rather than as a hidden failure.
This manuscript should therefore be read as a closed calibration and handoff package,
not as unfinished evidence waiting to be upgraded inside Energy itself.

\section{Reproducibility}
\label{sec:repro}

Current manuscript-side artifacts:
\begin{itemize}
  \item \texttt{stats/evidence/energy\_e1\_delta\_u\_table\_v1\_1.csv}
  \item \texttt{stats/evidence/energy\_e1\_tau\_summary\_v1\_1.csv}
  \item \texttt{stats/evidence/energy\_e2\_sensitivity\_summary\_v1.csv}
  \item \texttt{stats/evidence/recompute\_energy\_e1\_tables.py}
  \item \texttt{stats/evidence/recompute\_energy\_e2\_sensitivity.py}
  \item \texttt{scripts/rebuild\_all.sh}
  \item \path{results/analysis/energy_stagea_v2_impltrial_anthropic_gemini_summary_v1_2026-03-12.md}
  \item \path{results/analysis/energy_forward_main_batch_results_v1_2026-03-12.md}
  \item \path{results/analysis/energy_forward_main_batch_gemini_results_v1_2026-03-12.md}
  \item \path{results/analysis/energy_forward_combo_two_direction_anthropic_results_v1_2026-03-13.md}
  \item \path{results/analysis/energy_forward_combo_two_direction_gemini_results_v1_2026-03-13.md}
  \item \path{results/analysis/energy_priority_resolution_branch_handoff_v3_2026-03-18.md}
  \item \path{results/analysis/energy_priority_resolution_branch_synthesis_memo_v3_2026-03-18.md}
  \item \path{results/analysis/energy_priority_resolution_package_map_v1_2026-03-18.md}
  \item \path{results/analysis/energy_priority_resolution_manuscript_boundary_text_v2_2026-03-18.md}
  \item \path{results/analysis/energy_priority_resolution_package_boundary_external_review_decision_memo_v1_2026-03-18.md}
  \item \path{results/analysis/energy_policy_verification_side_topic_quarantine_readout_external_review_decision_memo_v1_2026-03-18.md}
  \item \path{results/analysis/energy_priority_resolution_horizon_stability_family_synthesis_memo_v1_2026-03-17.md}
\end{itemize}

For the current draft state, vendored calibration inputs are stored under:
\begin{itemize}
  \item \texttt{stats/stageb\_all/}
  \item \texttt{stats/combo\_rep3\_all/}
  \item \texttt{stats/raw\_stageb\_all/}
  \item \texttt{stats/raw\_combo\_rep3\_all/}
\end{itemize}

These inputs function here as frozen carried-forward paired inputs for the current Energy calibration surface. Detailed historical source and copy provenance are recorded in the reproducibility note rather than treated as part of the manuscript-facing role explanation.

\begin{thebibliography}{99}
\raggedright

\bibitem{saito2025nrr}
K.~Saito,
\newblock ``NRR-Core: Non-Resolution Reasoning as a Computational Framework for Contextual Identity and Ambiguity Preservation,''
\newblock \emph{arXiv preprint arXiv:2512.13478}, 2025.

\bibitem{saito2026phi}
K.~Saito,
\newblock ``NRR-Phi: Text-to-State Mapping for Ambiguity Preservation in LLM Inference,''
\newblock \emph{arXiv preprint arXiv:2601.19933}, 2026.

\bibitem{saito2026ime}
K.~Saito,
\newblock ``NRR-IME: Interface Decomposition for Stable State Transitions on Stateless LLM APIs,''
\newblock Repository preprint snapshot. \url{https://github.com/kei-saito-research/nrr-ime} (accessed 2026-03-05).

\bibitem{saito2026transfer}
K.~Saito,
\newblock ``NRR-Transfer: Cross-Domain Transfer of Phase 1.5 Operators Under Fixed Interface Constraints,''
\newblock Repository preprint snapshot. \url{https://github.com/kei-saito-research/nrr-transfer} (accessed 2026-03-05).

\bibitem{saito2026coupled}
K.~Saito,
\newblock ``NRR-Coupled (cp): Dependency-Consistency Evaluation for Coupled State Updates,''
\newblock Manuscript in preparation, 2026.

\bibitem{saito2026principles}
K.~Saito,
\newblock ``Integrated paper7 supporting package (comparison-led delayed-commitment surface),''
\newblock Repository preprint snapshot. \url{https://github.com/kei-saito-research/nrr-principles} (accessed 2026-03-07).

\bibitem{wilson1927}
E.~B. Wilson,
\newblock ``Probable inference, the law of succession, and statistical inference,''
\newblock \emph{Journal of the American Statistical Association}, 22(158):209--212, 1927.

\end{thebibliography}

\end{document}
